\documentclass[11pt]{article}
\usepackage[margin=1in]{geometry}
\usepackage{amsmath,amssymb,amsthm,mathtools}
\usepackage{enumitem}
\usepackage{hyperref}
\hypersetup{colorlinks=true,linkcolor=blue,citecolor=blue,urlcolor=blue}

\newtheorem{theorem}{Theorem}
\newtheorem{proposition}{Proposition}
\newtheorem{lemma}{Lemma}
\newtheorem{corollary}{Corollary}
\theoremstyle{definition}
\newtheorem{definition}{Definition}
\newtheorem{principle}{Principle}
\theoremstyle{remark}
\newtheorem{remark}{Remark}

\newcommand{\eps}{\varepsilon}
\newcommand{\C}{\mathbb{C}}
\newcommand{\PP}{\mathbb{P}^1}
\newcommand{\Stab}{\mathrm{Z}}
\newcommand{\diag}{\mathrm{diag}}
\newcommand{\Scat}{\mathcal{S}}
\newcommand{\Pmid}{P_{m\to m}}
\renewcommand{\Re}{\operatorname{Re}}
\newcommand{\lo}{\mathrm{lo}}
\newcommand{\mi}{\mathrm{mid}}
\newcommand{\hi}{\mathrm{hi}}
\newcommand{\Liouv}{\mathrm{Liouv}}

\title{\vspace{-2em}\textbf{Algebraic skeleton and transcendental Stokes core \\ of the
Type-1, $N{=}3$ Landau--Zener transition matrix}}
\author{Co-theoretical-physicist working session\\[2pt]
\small\texttt{projects/type-1-lz}; branch \texttt{claude/comathematician-agent-skills-9HPSJ}}
\date{2026-06-10}

\begin{document}
\maketitle
\vspace{-1.5em}

\begin{abstract}
The transition matrix of the Type-1, $N{=}3$ multistate Landau--Zener (MLZ) model is the
Stokes data of a meromorphic connection on $\PP$ with a single irregular singularity at
$u=\infty$ of Poincar\'e rank $2$. We organise its structure around one explicit
factorization and keep two mechanisms strictly apart. \emph{Architecture (wildness).} The
connection is wild (Theorem~\ref{thm:wild}); in the slope-ordered formal frame the
scattering matrix factors as $\Scat = U\,\Delta\,L$ along the two Stokes rays crossed by
the real axis, with $U,L$ unipotent and the diagonal $\Delta$ \emph{purely elementary}
(Theorem~\ref{thm:factor}). This pins the two extreme-slope survivals to the
Brundobler--Elser values and locates the only hard quantities, $\sigma$ and $b$, in
disjoint unipotent multiplier slots (Proposition~\ref{prop:slots}). \emph{Hardness
(transcendence).} That these slot values are non-elementary is a consequence of
\textbf{non-rigidity}, not of wildness: $N{=}2$ Landau--Zener is wild yet has an elementary
connection constant, so wildness cannot be the source (\S\ref{sec:deriv}). We derive the
true source: a rigid wild connection has a $\Gamma$-ratio constant (Katz/middle
convolution), whereas the non-rigid Type-1 connection ($\chi=0$, one accessory parameter;
$\dim=6$ wild character variety) has a connection constant that is an isomonodromy
(Painlev\'e/Garnier) transcendent, hence non-Liouvillian (Theorem~\ref{thm:transc}). A
differential-algebraic order test confirms in-model that $\sigma(\text{parameter})$ is
\emph{not} order-$1$ differentially algebraic while behaving as an order-$2$ transcendent
(\S\ref{sec:cert}). \textbf{Honest scope.} The wildness, factorization, slot, non-rigidity
and empty-cell statements are rigorous or rigorous-plus-numerically-verified; the final
\emph{non-Liouvillian} step is complete for the Painlev\'e VI reduction and a flagged
frontier item for the full rank-$3$ Garnier flow. This document supersedes the earlier note
\texttt{type1\_n3\_wildness\_proof.tex}, correcting one claim (the formal monodromy is
diagonal, not a $3$-cycle; Remark~\ref{rem:retract}).
\end{abstract}

\tableofcontents

\section{The model and the objects}\label{sec:model}

\paragraph{Physics first.} A quantum system with three levels is swept once through an
avoided crossing: its Hamiltonian is $H(u)=H_0+uA$, where $u$ runs from $-\infty$ to
$+\infty$ (proportional to time), $A=\diag(a_0,a_1,a_2)$ holds the constant diabatic
\emph{slopes}, and $H_0$ holds the couplings. We read off the survival and transition
probabilities $P_{x\to j}$ --- the probability of entering on diabatic level $x$ and exiting
on level $j$. The matrix $P$ is doubly stochastic. The single quantity this paper is built
around is the \emph{middle-level survival} $\sigma$: the probability that a system entering
on the middle slope leaves on it.

\paragraph{Type-1 data.} Fix $\eps_0<\eps_1<\eps_2$, real couplings $\gamma_i$, and distinct
slopes $a_0,a_1,a_2$. The Type-1 $N{=}3$ MLZ Hamiltonian is $H(u)=H_0+uA$ with the
Cauchy/Gaudin couplings
\begin{equation}\label{eq:H0}
(H_0)_{ij}=\frac{\gamma_i\gamma_j\,(a_i-a_j)}{\eps_i-\eps_j}\ (i\neq j),\qquad
(H_0)_{ii}=-\sum_{k\neq i}\frac{\gamma_k^2\,(a_i-a_k)}{\eps_i-\eps_k}.
\end{equation}
Type-1 is the integrable structure isolated in the classification of parameter-dependent
quantum integrable models \cite{OWY2009,OY2011}: the family
$\{H_0(a)+u\,\diag(a)\}$ is a commuting (Gaudin-magnet) family sharing one eigenbasis for
all $u$ \cite{Y2018}, and \eqref{eq:H0} is the solvable MLZ representative
\cite{SC2017}. The Schr\"odinger system $i\,\psi'(u)=H(u)\psi(u)$ defines a meromorphic
connection on the trivial rank-$3$ bundle over $\PP_u$,
\begin{equation}\label{eq:conn}
\nabla \;=\; \frac{d}{du}\;-\;i\big(H_0+uA\big),
\end{equation}
holomorphic on $\C_u$ and singular only at $u=\infty$.

\paragraph{The observable.} The scattering matrix $\Scat$ relates the canonical asymptotic
frames of \eqref{eq:conn} at $u=-\infty$ and $u=+\infty$; the transition matrix is
$P_{x\to j}=|\Scat_{xj}|^2$, doubly stochastic (indeed unistochastic). Order the channels by
slope, $\lo<\mi<\hi$ (the slope-ascending order of $a$), and write
$\Pmid=|\Scat_{mm}|^2$ for the middle survival.

\paragraph{Geometric dictionary (for the reader from the other side).} Equation
\eqref{eq:conn} is a rank-$3$ meromorphic connection with one singular point. ``Wild''
means that point is irregular (the formal solutions carry exponential factors); the
\emph{exponential torus}, \emph{formal monodromy} and \emph{Stokes matrices} are its local
formal-plus-analytic invariants \cite{Wasow1965,LR2016,vdPS2003}; the \emph{wild character
variety} is the moduli of such connections with fixed formal type \cite{Boalch2014}.
``Rigid'' means the global Stokes data are determined by the local data
(Katz index $2$) \cite{Katz1996}. Each notion is defined where first used.

\paragraph{The two hard numbers, and the elementary skeleton around them.} $P$ has four real
degrees of freedom. Two are fixed exactly and elementarily by the \textbf{Brundobler--Elser
(BE) law} \cite{BE1993} (proved in \cite{VO2004,DS2006}): the two extreme-slope survivals
are
\begin{equation}\label{eq:BE}
P_{\lo\to\lo}=e^{-2\pi(\delta_{lm}+\delta_{lh})},\quad
P_{\hi\to\hi}=e^{-2\pi(\delta_{lh}+\delta_{mh})},\qquad
\delta_{ij}=\frac{\gamma_i^2\gamma_j^2\,|a_i-a_j|}{(\eps_i-\eps_j)^2}.
\end{equation}
Imposing \eqref{eq:BE} together with double stochasticity leaves $P$ determined by exactly
two further numbers, which we take to be
\begin{equation}\label{eq:sigmab}
\sigma:=\Pmid \qquad\text{and}\qquad b:=P_{\hi\to\lo}.
\end{equation}
We treat $\sigma$ and $b$ symmetrically: both are genuine \emph{Stokes constants} of the
wild connection (Corollary~\ref{cor:stokes}), and the wild core they span is two-dimensional
(\S\ref{sec:cert}). They are not on the same footing physically --- $\sigma$ is hard in
every regime, whereas $b$ is elementary in the diabatic and adiabatic limits --- but neither
is reducible to the other (Proposition~\ref{prop:slots}, \S\ref{sec:cert}). Canonical anchor,
used throughout for numerics: $\eps=(-2,0,3)$, $\gamma=(1,0.8,1.2)$, $a=(-1,0.5,2)$, for
which $\sigma=0.214724\ldots$

\begin{definition}[Wild/tame; rigid]\label{def:terms}
A meromorphic connection is \emph{tame} (regular-singular) at a point if every formal
solution there grows at most polynomially; otherwise it is \emph{wild}. For an irreducible
connection on $\PP$ with singular set $S$ and local data $g_s$, the \emph{Katz rigidity
index} is $\chi=(2-|S|)N^2+\sum_{s\in S}\dim \Stab(g_s)$, $\Stab$ the centralizer in
$GL_N(\C)$; the connection is \emph{rigid} iff $\chi=2$, equivalently its global
monodromy/Stokes data are determined by the local data \cite{Katz1996}. A function of the
model parameters is \emph{Liouvillian} if it lies in a tower over $\C(\eps,\gamma,a)$ built
from algebraic extensions, exponentials and integrals (equivalently, has solvable
differential Galois group \cite{Kolchin1973,vdPS2003}).
\end{definition}

\section{Wildness: the architecture}\label{sec:wild}

\begin{theorem}[Wildness]\label{thm:wild}
The connection \eqref{eq:conn} has at $u=\infty$ an irregular singularity of Poincar\'e
rank $2$. It is \emph{unramified}, with $N{=}3$ distinct exponential rates
$q_i(u)=\tfrac{i}{2}a_i u^2$, and its Stokes structure is non-trivial. In particular
$\nabla$ is wild, not Fuchsian.
\end{theorem}

\begin{proof}
$M(u)=i(H_0+uA)$ is degree $1$ in $u$. In $t=1/u$ at $\infty$, $\frac{d}{du}=-t^2\frac{d}{dt}$,
so $\nabla$ becomes $\frac{d}{dt}+i\big(A/t^{3}+H_0/t^{2}\big)$, a pole of order $3$;
Poincar\'e rank $=$ (pole order)$-1=2$. The leading coefficient $iA/t^3$ is regular semisimple
(the $a_i$ are distinct), so the point is unramified: the formal solution diagonalizes with
exponential factors $\exp\!\big(\int^{u} i a_i u'\,du'\big)=\exp(\tfrac{i}{2}a_i u^2)$, i.e.
distinct quadratic rates $q_i$ \cite{Wasow1965}. The rates are non-constant, so no gauge
removes the exponentials and $\nabla$ is wild. For each pair $i\neq j$ the dominance of
$e^{q_i}$ over $e^{q_j}$ switches across the rays $\Re[(a_i-a_j)u^2]=0$, and the associated
Stokes matrices are non-trivial (for $N{=}2$ this is Weber's equation, whose Stokes
multiplier is the elementary Landau--Zener factor \cite{JMU1981}).
\end{proof}

\begin{corollary}\label{cor:stokes}
$\Scat$, and hence both $\sigma$ and $b$, are \emph{Stokes connection constants} of the wild
connection \eqref{eq:conn}: points of its wild character variety, fixed by the global gluing
of the Stokes sectors at $\infty$, not by any Fuchsian monodromy.
\end{corollary}

\section{The explicit construction: $\Scat=U\,\Delta\,L$}\label{sec:constr}

This section is the centerpiece. We make the architecture of Theorem~\ref{thm:wild} explicit:
the Stokes rays, the factorization of $\Scat$ across them, the elementary diagonal, and the
exact location of $\sigma,b$ in the unipotent multiplier slots. Heavier analytic details are
in Appendix~\ref{app:sectorial}.

\begin{lemma}[Maximal ray-degeneracy]\label{lem:rays}
All three pairwise rate-differences are pure-imaginary quadratics,
$q_i-q_j=-\tfrac{i}{2}(a_i-a_j)u^2$ (the $a$ are real). On $u=Re^{i\theta}$,
\[
\Re\big(q_i-q_j\big)=\tfrac12(a_i-a_j)R^2\sin 2\theta,
\]
a single function of $\theta$ for \emph{all} pairs. Consequently the three pairs share the
same four \emph{anti-Stokes} directions (the axes $\theta=0,\tfrac\pi2,\pi,\tfrac{3\pi}2$)
and the same four \emph{Stokes/jump} rays (the diagonals $\theta=\pm\tfrac\pi4,\pm\tfrac{3\pi}4$);
the dominance order of $\{e^{q_j}\}$ equals the slope order in the quadrants $1,3$ and its
reverse in $2,4$. The physical scattering directions $\theta=\pi,0$ are anti-Stokes, which is
why $P=|\Scat|^2$ is unistochastic.
\end{lemma}

\begin{proof}
Immediate from $a_i\in\mathbb{R}$ and $\Re(-\tfrac{i}{2}(a_i-a_j)R^2 e^{2i\theta})
=\tfrac12(a_i-a_j)R^2\sin2\theta$. Generic irregular data of rank $2$ would give $12$
distinct rays for the three pairs; reality of the commuting slopes collapses them to $4$.
\end{proof}

\begin{theorem}[Lateral factorization with elementary diagonal]\label{thm:factor}
In the slope-ordered formal frame at $u=\infty$, the scattering matrix factors as
\begin{equation}\label{eq:UDL}
\Scat \;=\; U\,\Delta\,L,
\end{equation}
where $U$ is unit upper-triangular (the Stokes factor on the ray $\theta=\tfrac\pi4$), $L$ is
unit lower-triangular (the ray $\theta=\tfrac{3\pi}4$), and $\Delta=\diag(\Delta_\lo,\Delta_\mi,\Delta_\hi)$
is the formal transport through the sector containing $\theta=\tfrac\pi2$. The diagonal is
\emph{purely elementary}:
\begin{equation}\label{eq:Delta}
|\Delta_\mi|=e^{-\pi|b_\mi|},\quad |\Delta_\hi|=e^{-\pi|b_\hi|},\quad
|\Delta_\lo|=e^{+\pi|b_\lo|},\qquad
b_j:=\sum_{k\neq j}\frac{(H_0)_{jk}^2}{a_j-a_k},
\end{equation}
the $b_j$ being the signed BE exponents ($\sum_j b_j=0$).
\end{theorem}

\begin{proof}[Proof (structure; analytic details Appendix~\ref{app:sectorial})]
By Lemma~\ref{lem:rays} the real axis from $\theta=\pi$ to $\theta=0$ through the upper
half-plane crosses exactly the two jump rays $\tfrac{3\pi}4$ then $\tfrac\pi4$; at each, all
three pairs are simultaneously dominant/recessive (maximal degeneracy), so each carries a
\emph{full} unipotent Stokes factor, the two oppositely triangular in slope order. The
unramified formal solution is $\hat F(u)\,u^{L_0}e^{Q(u)}$, $Q=\diag(q_i)$, with $\hat F$ a
formal gauge \cite{Wasow1965,LR2016}; the diagonal transport through the intervening sector
is $\Delta$, whose moduli are the half-turn moduli $e^{\mp\pi b_j}$ of the formal monodromy.
The mid and hi entries are the recessive/standard BE weights $e^{-\pi|b_j|}$; the lo entry is
the dual weight $e^{+\pi|b_\lo|}$, forced by $|\det\Scat|=1=\prod_j|\Delta_j|$ together with
$\sum_jb_j=0$ (the last Gauss pivot absorbs the unitarity normalization). Numerically,
\eqref{eq:UDL}--\eqref{eq:Delta} hold to the integration floor across all sampled data: the
extracted $|\Delta_j|/e^{\mp\pi|b_j|}$ equal $(1,1,1)$ to $\le 10^{-3}$ (three cases, two
gold-gated; \S\ref{sec:cert}, Appendix~\ref{app:sectorial}), and the lower-triangular order
(LDU) does \emph{not} have an elementary diagonal, which is what selects $U\Delta L$ as the
upper-half-plane lateral product.
\end{proof}

\begin{proposition}[Slot equations]\label{prop:slots}
Write $\Scat=U\Delta L$ in slope order $(0,1,2)=(\lo,\mi,\hi)$. Then
\begin{align}
P_{\hi\to\hi}&=|\Delta_\hi|^2 &&\text{(BE, manifest, multiplier-free)},\label{eq:Phh}\\
P_{\lo\to\lo}&=\big|\Delta_\lo+U_{01}\Delta_\mi L_{10}+U_{02}\Delta_\hi L_{20}\big|^2=e^{-2\pi|b_\lo|}
&&\text{(BE, as a forced cancellation)},\label{eq:Pll}\\
b=P_{\hi\to\lo}&=|L_{20}|^2\,|\Delta_\hi|^2 &&\text{(\emph{one} multiplier)},\label{eq:b}\\
\sigma=P_{\mi\to\mi}&=\big|\Delta_\mi+U_{12}\Delta_\hi L_{21}\big|^2 &&\text{(\emph{two} multipliers + a phase).}\label{eq:sigma}
\end{align}
The remaining entries are fixed affinely by double stochasticity. Thus $b$ depends only on the
slot $L_{20}$, while $\sigma$ depends only on $\{U_{12},L_{21}\}$ and their relative phase:
the two hard numbers occupy \emph{disjoint} multiplier slots.
\end{proposition}

\begin{proof}
Direct expansion of $|(U\Delta L)_{xj}|^2$ in slope order; \eqref{eq:Phh} because $\Scat_{\hi\hi}=\Delta_\hi$
(no slot reaches the top diagonal in $U\Delta L$), \eqref{eq:Pll} because the bottom diagonal
collects all three lateral contributions and equals the BE weight by Theorem~\ref{thm:factor},
\eqref{eq:b} from $\Scat_{\hi\lo}=L_{20}\Delta_\hi$, and \eqref{eq:sigma} from
$\Scat_{\mi\mi}=\Delta_\mi+U_{12}\Delta_\hi L_{21}$. All four are exact identities at solver
precision in every sampled case (\S\ref{sec:cert}).
\end{proof}

Read physically, \eqref{eq:sigma} says $\sigma$ is the interference of the direct formal
mid-passage (amplitude the naive BE weight $e^{-\pi|b_\mi|}$) with the single recombination
path $\mi\!\to\!\hi\!\to\!\mi$ through the two ray crossings; the deviation of $\sigma$ from
its BE value \emph{is} that recombination term, as an equation rather than a metaphor.

\begin{remark}[Two factorizations, one skeleton --- only one is exact]\label{rem:contrast}
The same extreme survivals and oriented cycle are reproduced by the elementary
\emph{time-ordered} product $M_3M_2M_1$ of three $2$-level Landau--Zener stochastic factors,
composed in the order of the rational diabatic crossing times
$u_{ij}=\big((H_0)_{jj}-(H_0)_{ii}\big)/(a_i-a_j)$ (the incoherent skeleton). This product
is exact on $P_{\lo\to\lo},P_{\hi\to\hi}$ and the cycle, but returns \emph{wrong} values for
$\sigma$ and $b$ (off by up to $\sim0.99$ for $\sigma$ in deep overlap). It composes in
\emph{$u$-time order}; the exact $U\Delta L$ composes in \emph{ray order at $\infty$}. The
gap between them is precisely the Stokes core $\{\sigma,b\}$.
\end{remark}

\begin{remark}[Correction to the predecessor note]\label{rem:retract}
An earlier version (\texttt{type1\_n3\_wildness\_proof.tex}) stated that the formal monodromy,
read on $P$, is an oriented $3$-cycle. This is incorrect: at an \emph{unramified} irregular
point the formal monodromy is \emph{diagonal}, and its half-turn moduli $e^{\mp\pi b_j}$ are
exactly the BE weights $\Delta$ of Theorem~\ref{thm:factor}. The oriented $3$-cycle is real
but belongs to the incoherent skeleton of Remark~\ref{rem:contrast} (its deep-overlap limit);
its orientation is the $\mathbb{Z}_2$ selector $\mathrm{sgn}(\pi)$, the parity of the
slope-ordering permutation \cite{OWY2009}. We assign each object its correct role here.
\end{remark}

\section{Non-rigidity, and the true source of transcendence}\label{sec:deriv}

Section~\ref{sec:constr} put $\sigma,b$ in explicit elementary slots. That the slot
\emph{values} are non-elementary is a separate fact, with a separate cause. We first record
non-rigidity, then derive --- as a mechanism, not an analogy --- that non-rigidity (and not
wildness) is what makes the slot values transcendental.

\subsection{Non-rigidity}

\begin{lemma}[Tame de-confluence]\label{lem:deconf}
Replacing the unbounded ramp $uA$ by a bounded saturating ramp $f(u)A$, $f(\pm\infty)=\pm1$
(e.g.\ $f=\tanh$), and setting $z=\tfrac12(1+f(u))$, yields a rank-$3$ Fuchsian system on
$\PP$ with three regular singular points $z\in\{0,1,\infty\}$ and residues
$R_0=\tfrac{1}{2i}(H_0-A)$, $R_1=-\tfrac{1}{2i}(H_0+A)$, $R_\infty=-iA$, all regular
semisimple for generic Type-1 data. The wild connection \eqref{eq:conn} is the confluent
limit in which $z=0,1$ collide at $\infty$ into the rank-$2$ irregular point.
\end{lemma}

\begin{proof}
With $z=\tfrac12(1+\tanh u)$, $dz=2z(1-z)\,du$ gives
$\partial_z\psi=\frac{H_0+(2z-1)A}{2i\,z(1-z)}\psi$, simple poles at $z=0,1$ with the stated
residues and (from the $A/(-iz)$ decay) at $z=\infty$ with $R_\infty=-iA$; eigenvalues
$-ia_k$ and $\tfrac1{2i}\mathrm{spec}(H_0\mp A)$ are distinct generically. Confluence is the
standard Heun$\to$rank-$2$ degeneration \cite{Ronveaux1995,JMU1981}.
\end{proof}

\begin{theorem}[Non-rigidity]\label{thm:nonrigid}
The connection \eqref{eq:conn} is non-rigid: its wild character variety has positive
dimension and the Stokes data are not determined by the local formal data.
\end{theorem}

\begin{proof}
By Lemma~\ref{lem:deconf} the de-confluenced system is rank $3$, three points, all residues
regular semisimple, $\dim\Stab=3$, so the Katz index is
$\chi=(2-3)\cdot3^2+3\cdot3=0<2$; it is non-rigid with $\tfrac12(2-\chi)=1$ accessory
parameter (Heun/Garnier type) \cite{Katz1996}. Confluence with fixed formal type does not
lower the moduli dimension \cite{JMU1981,Boalch2014}, so \eqref{eq:conn} is non-rigid as well.
Independently, the Boalch dimension formula for the wild character variety at $(N,r)=(3,2)$
with distinct leading rates gives $\dim=6>0$ \cite{Boalch2014}; the single Heun accessory
parameter is the leading modulus, the remaining dimensions being the framing/Stokes torus.
\end{proof}

\subsection{The derivation: wildness cannot be the cause; non-rigidity is}

\begin{lemma}[The $N{=}2$ separation]\label{lem:N2}
$N{=}2$ Landau--Zener is wild, with non-Liouvillian fundamental solutions in $u$
(parabolic-cylinder functions), yet its connection constant is elementary,
$P_{1\to1}=e^{-2\pi\delta}$. Hence ``no closed-form solution in $u$'' (a statement about the
differential Galois group in $u$) does \emph{not} imply ``transcendental connection
constant'': wildness alone cannot make a Stokes constant non-elementary.
\end{lemma}

\begin{proof}
Weber's equation is the rank-$2$ unramified case of Theorem~\ref{thm:wild}; its exponential
torus and two opposite Stokes unipotents generate a non-solvable subgroup of $GL_2$, so by
the Ramis density theorem and Kolchin's criterion the solutions are non-Liouvillian in $u$
\cite{Ramis1985,vdPS2003,Kolchin1973}. The single Stokes multiplier nonetheless evaluates to
the elementary factor $e^{-2\pi\delta}$ \cite{JMU1981}. The two statements are independent.
\end{proof}

What distinguishes $N{=}2$ (elementary constant) from Type-1 $N{=}3$ (hard constant) is not
wildness --- both are wild --- but \emph{rigidity}.

\begin{proposition}[Rigid $\Rightarrow$ elementary constant]\label{prop:rigid}
If an irreducible (wild) connection is rigid, its connection/Stokes data are explicit
$\Gamma$-ratios in the local exponents.
\end{proposition}

\begin{proof}
A rigid connection is built from rank-$1$ objects by middle convolution and tensor twists
(Katz; algorithmic form Dettweiler--Reiter) \cite{Katz1996,DR2000,DR2007}; each operation is
an explicit Euler/Laplace integral, so the connection matrix is a finite product of
$\Gamma$-factors of the local exponents. $N{=}2$ LZ realises this, returning $e^{-2\pi\delta}$.
\end{proof}

\begin{theorem}[Non-rigid $\Rightarrow$ Painlev\'e/Garnier transcendent $\Rightarrow$ non-Liouvillian]\label{thm:transc}
For generic Type-1 $N{=}3$ data the connection constants $\sigma,b$ are non-Liouvillian
functions of $(\eps,\gamma,a)$.
\end{theorem}

\begin{proof}[Derivation (with the rigor ceiling stated)]
\emph{Step 1 (deformation $\Leftrightarrow$ non-rigidity).} By Theorem~\ref{thm:nonrigid} the
connection is a non-isolated point of a $6$-dimensional moduli space carrying one accessory
parameter; $\sigma,b$ are non-constant coordinates on it.
\emph{Step 2 (the deformation is an isomonodromy/Painlev\'e flow).} Realise the accessory
parameter as the position $\lambda$ of an apparent singularity of the de-confluenced system
(Lemma~\ref{lem:deconf}); isomonodromic motion of $\lambda$ obeys the Schlesinger/JMU
equations \cite{JMU1981}, of Garnier type here (the rank-$3$, three-point analog of
Heun$\leftrightarrow$Painlev\'e\,VI). In the wild picture these are Boalch's wild isomonodromy
flows on the character variety, which realise Painlev\'e/Garnier systems geometrically
\cite{Boalch2001,Boalch2007}. Along the flow the connection constant solves a second-order
nonlinear (Painlev\'e/Garnier) equation, not a linear one.
\emph{Step 3 (irreducibility $\Rightarrow$ non-Liouvillian).} The general solution of such an
equation lies in no Liouvillian extension of the coefficient field: for Painlev\'e I--VI this
is the Nishioka--Umemura irreducibility theorem \cite{Nishioka1988,Umemura1990} and the
Malgrange-groupoid proofs of Casale \cite{Casale2008}; for Painlev\'e VI it is proved by
Cantat--Loray via the dynamics on the character variety \cite{CL2009} --- exactly our setting,
since the wild character variety \emph{is} that variety and the flow \emph{is} the dynamics
on it. Pulling back along the algebraic embedding of Type-1 data, $\sigma,b$ are
non-Liouvillian over $\C(\eps,\gamma,a)$; the appropriate frame is parameterized
Picard--Vessiot theory over a differential parameter field \cite{CS2007}.
\emph{Step 4 (genericity).} This holds off the classical-solution locus (algebraic/Riccati
special solutions, codimension $\ge1$). Type-1 data is non-classical, witnessed in-model by
the two-dimensional transcendence degree and the order test of \S\ref{sec:cert}.

\emph{Rigor ceiling.} Steps $1$--$2$ and the $N{=}2$ separation are rigorous; Step $3$ is
complete for Painlev\'e I--VI (including the PVI character-variety-dynamics proof matching our
frame), but for the genuine rank-$3$ \emph{Garnier} flow the irreducibility is established
only in major cases. This single implication is the honest frontier item; everything else in
the theorem is rigorous or numerically certified.
\end{proof}

\subsection{The disaggregation made into a principle}

\begin{principle}[Architecture vs hardness]\label{prin:disagg}
Two independent axes govern $P$. \emph{Wildness} (irregular type $A=\diag(a)$) supplies the
\emph{architecture}: the Stokes rays, the factorization $U\Delta L$, the elementary $\Delta$,
the BE-pinning, and the disjoint slots holding $\sigma,b$. \emph{Non-rigidity} (positive-%
dimensional moduli) supplies the \emph{hardness}: it makes the slot \emph{values}
Painlev\'e/Garnier transcendents, hence non-Liouvillian. The axes are independent:
\[
\begin{array}{c|cc}
 & \text{rigid} & \text{non-rigid}\\\hline
\text{wild} & N{=}2\ \text{LZ: elementary }e^{-2\pi\delta} & \boxed{\text{Type-1 }N{=}3:\ \sigma,b\ \text{transcendental}}\\
\text{tame} & {}_3F_2/\Gamma\text{-ratios} & \text{Heun/PVI: transcendental connection constant}
\end{array}
\]
$N{=}2$ (wild, rigid) is elementary; Heun (tame, non-rigid) is transcendental; Type-1 $N{=}3$
is wild \emph{and} non-rigid, and its transcendence rides on the non-rigid axis only. Heun is
thus the base case of the \emph{same} mechanism (Theorem~\ref{thm:transc}), not an analogy.
\end{principle}

\section{The empty cell}\label{sec:empty}

\begin{theorem}[No closed form by any rigidity mechanism]\label{thm:empty}
In the $(\text{tame/wild})\times(\text{rigid/non-rigid})$ classification, the Type-1 $N{=}3$
amplitude occupies the cell \emph{wild $\wedge$ non-rigid}, which admits no closed-form
mechanism: it is neither a tame (Fuchsian) monodromy datum nor a rigid (generalized-%
hypergeometric/$\Gamma$-function) connection constant. Its hard data is the
two-dimensional Stokes core $\{\sigma,b\}$ of Principle~\ref{prin:disagg}.
\end{theorem}

\begin{proof}
Wild by Theorem~\ref{thm:wild}; non-rigid by Theorem~\ref{thm:nonrigid}. Tameness is
contradicted by the exponential rates; the rigid-mechanism closed form
(Proposition~\ref{prop:rigid}) requires $\chi=2$, contradicted by $\chi=0$. The three
solvable neighbors (tame-rigid ${}_3F_2$; wild-rigid $N{=}2$ LZ and solvable MLZ; tame-non-%
rigid Heun) each fail one hypothesis. Hence no rigidity mechanism applies, and the residual
data is $\{\sigma,b\}$ (\S\ref{sec:cert}).
\end{proof}

\section{Numerical certification of the core}\label{sec:cert}

We certify, gold-gated against $\sigma=0.214724$, the two quantitative claims that the
analytic statements rest on: the core is two-dimensional, and $\sigma$ is a transcendent (not
order-$1$ differentially algebraic). Scripts: \texttt{ws\_stokes\_slots2.py} (factorization,
slots), \texttt{ws\_transcendence\_degree.py} (trdeg), \texttt{ws\_painleve\_da\_test.py}
(order test).

\begin{proposition}[The Stokes core is two-dimensional]\label{prop:trdeg}
$\{\sigma,b\}$ has transcendence degree two over the elementary field: $b$ is not a function
of $(P_{\lo\to\lo},P_{\hi\to\hi},\sigma)$.
\end{proposition}

\begin{proof}[Evidence (numerically-supported)]
The Jacobian of $(\eps,\gamma,a)\mapsto(P_{\lo\to\lo},P_{\hi\to\hi},\sigma,b)$ attains full
rank $4$ across the family (e.g.\ singular values $(1.50,0.39,0.35,0.21)$ at a random anchor,
the smallest $\sim200\times$ the floor), so $\Phi$ is a submersion and $b$ is not pinned by
the triple. The best polynomial fit of $b$ on $(P_{\lo\to\lo},P_{\hi\to\hi},\sigma)$ leaves an
irreducible residual $\approx0.26$ against $\mathrm{std}(b)\approx0.35$, far above the
$10^{-3}$ floor. The disjoint-slot structure \eqref{eq:b}--\eqref{eq:sigma} is the analytic
shadow of this independence.
\end{proof}

\begin{proposition}[$\sigma$ is a transcendent, not an elementary function]\label{prop:da}
Along a coupling-scale slice $\gamma\mapsto\lambda\gamma$ on which $\sigma$ sweeps
$[0.016,0.905]$, $\sigma(\lambda)$ is \emph{not} order-$1$ differentially algebraic, while
behaving as an order-$2$ transcendent.
\end{proposition}

\begin{proof}[Evidence (numerically-supported; order-$2$ structure suggestive)]
Computing $\sigma(\lambda)$ to $\sim10^{-6}$, a Chebyshev interpolant and a search for an
algebraic relation $P(\lambda,\sigma,\dots,\sigma^{(k)})=0$ (smallest singular value of the
normalized monomial-jet matrix) gives, at the only well-conditioned cell (order-$1$,
degree-$2$): an elementary control $e^{-c(\lambda)^2}$ a clean hit ($4\times10^{-15}$), a
holonomic control $J_0$ a clean miss ($3\times10^{-3}$) with an order-$2$ hit
($7\times10^{-12}$), and $\sigma$ a \emph{miss} at $1.5\times10^{-4}$ --- $\sim37\times$ its
floor, the same order as $J_0$'s miss --- with order-$2$ residual at floor. The miss is stable
under resolution ($1.50\times10^{-4}$ at degree $20$, $1.48\times10^{-4}$ at degree $26$ while
the floor improved $60\times$), hence a real ODE-obstruction, not fit error. Thus $\sigma$
tracks the transcendent pattern (order-$1$ no, order-$2$ yes), not the elementary pattern
(order-$1$ yes), consistent with Theorem~\ref{thm:transc}. Pinning the exact order-$2$
(Painlev\'e/Garnier) structure --- order-$2$ residual orders below floor, as $J_0$ achieves
--- needs $\sigma$ to $\sim10^{-9}$ (a high-precision connection solver), a flagged item.
\end{proof}

\section{Conclusion}\label{sec:concl}

The Type-1 $N{=}3$ transition matrix separates cleanly into an explicit \emph{algebraic
skeleton} and a \emph{transcendental Stokes core}, and the separation is causal, not
descriptive. Wildness --- the irregular type $A=\diag(a)$ --- builds the skeleton: it fixes
the four shared Stokes rays, the lateral factorization $\Scat=U\Delta L$, the elementary
diagonal $\Delta$ of BE weights, and the two disjoint multiplier slots in which $\sigma$ and
$b$ sit (Theorems~\ref{thm:wild},~\ref{thm:factor}, Proposition~\ref{prop:slots}).
Non-rigidity --- the positive-dimensional moduli --- supplies the only hardness: through
isomonodromy it makes the slot values Painlev\'e/Garnier transcendents, hence non-Liouvillian
(Theorem~\ref{thm:transc}), a property wildness alone cannot confer ($N{=}2$ is the witness).
``Essential wildness'' and ``transcendence'' are therefore two facts, not one: the first
locates the data, the second makes it hard. The core is two-dimensional, $\{\sigma,b\}$, and
neither number admits a closed form by any rigidity mechanism --- not because the connection
is wild, but because it is non-rigid.

\appendix

\section{Sectorial existence and the elementary diagonal}\label{app:sectorial}

We sketch the analytic backbone of Theorem~\ref{thm:factor}; full proofs follow the standard
sectorial theory \cite{Wasow1965,LR2016}. On each of the four sectors bounded by the jump
rays of Lemma~\ref{lem:rays} there is a unique sectorial fundamental solution asymptotic to
$\hat F(u)u^{L_0}e^{Q}$ (Sibuya/Hukuhara normal solutions). The Stokes matrices relating
consecutive sectors are unipotent, upper or lower triangular according to the dominance order,
which by Lemma~\ref{lem:rays} alternates by quadrant; crossing the upper half-plane from
$\theta=\pi$ to $0$ therefore composes one lower and one upper unipotent with the diagonal
formal transport, giving \eqref{eq:UDL}. The diagonal moduli are the half-turn factors of the
formal monodromy $u^{L_0}$, whose diagonal entries are $-i\,b_j$ with $b_j$ the signed BE
exponents; analytic continuation across the upper half-plane contributes $\log u\mapsto\log
u+i\pi$, giving $|\Delta_j|=e^{\mp\pi|b_j|}$ with the dual ($+$) sign forced on the last Gauss
pivot by $\det$. The elementary-diagonal claim is the assertion that $\hat F$ contributes no
modulus to the diagonal at the anti-Stokes endpoints --- a Schur-complement (minor) identity
for lateral products, which we verify numerically to $\le10^{-3}$ (Table below) and expect to
admit a closed proof.

\smallskip
\noindent\emph{Verification (\texttt{ws\_stokes\_slots2.py}, Richardson in the cutoff $R$):}
extracted $|\Delta_j|/e^{\mp\pi|b_j|}$ extrapolate to $(1,1,1)$ ---
canonical $(13.378,1.00004,1.00003)$ with $1/\mathrm{BE}_\lo^2=13.380$;
sampleB $(2482.9,0.9995,0.9996)$; a third case $(52.420,1.00002,0.99998)$ ---
and the slot equations \eqref{eq:Phh}--\eqref{eq:sigma} reproduce $\sigma$ to
$\le10^{-5}$ (gold $0.214724$) with $\|\Scat\Scat^\dagger-I\|\le10^{-5}$.

\section{The differential-algebraic order test}\label{app:da}

For a function $f$ of one parameter sampled to accuracy $\eta$, build a Chebyshev interpolant,
obtain $f',f''$ by spectral differentiation, and for each order $k$ and degree $d$ form the
matrix whose columns are the degree-$\le d$ monomials in $(t,f,f',\dots,f^{(k)})$ evaluated at
interior points; its smallest singular value (after column normalization) is small iff an
algebraic ODE of that $(k,d)$ holds. Only low $(k,d)$ are diagnostic: at high degree the
monomials become numerically dependent over a smooth curve (Vandermonde ill-conditioning) and
the singular value collapses spuriously --- visible in the controls (the holonomic $J_0$ shows
a spurious order-$1$ ``hit'' only at degree $\ge3$). The robust signature is the comparison at
order-$1$, degree-$2$: elementary (Liouvillian) functions hit; genuine transcendents
(holonomic or Painlev\'e) miss. The stability of the miss under increasing fit degree
distinguishes a true ODE-obstruction from interpolation error. See
\texttt{ws\_painleve\_da\_test.py} and Proposition~\ref{prop:da}.

\begin{thebibliography}{99}

\bibitem{BE1993} S.~Brundobler, V.~Elser, \emph{S-matrix for generalized Landau--Zener
problem}, J.\ Phys.\ A \textbf{26} (1993) 1211--1227.

\bibitem{SC2017} N.~A.~Sinitsyn, V.~Y.~Chernyak, \emph{The quest for solvable multistate
Landau--Zener models}, J.\ Phys.\ A \textbf{50} (2017) 255203, arXiv:1701.01870.

\bibitem{OWY2009} H.~K.~Owusu, K.~Wagh, E.~A.~Yuzbashyan, \emph{The link between integrability,
level crossings and exact solution in quantum models}, J.\ Phys.\ A \textbf{42} (2009)
035206, arXiv:0807.0259.

\bibitem{OY2011} H.~K.~Owusu, E.~A.~Yuzbashyan, \emph{Classification of parameter-dependent
quantum integrable models, their parameterization, exact solution and other properties},
J.\ Phys.\ A \textbf{44} (2011) 395302, arXiv:1106.1831.

\bibitem{Y2018} E.~A.~Yuzbashyan, \emph{Integrable time-dependent Hamiltonians, solvable
Landau--Zener models and Gaudin magnets}, Ann.\ Phys.\ \textbf{392} (2018) 323--339,
arXiv:1804.04926.

\bibitem{VO2004} M.~V.~Volkov, V.~N.~Ostrovsky, \emph{Exact results for survival probability
in the multistate Landau--Zener model}, J.\ Phys.\ B \textbf{37} (2004) 4069--4084.

\bibitem{DS2006} B.~E.~Dobrescu, N.~A.~Sinitsyn, \emph{Comment on `Exact results for survival
probability in the multistate Landau--Zener model'}, J.\ Phys.\ B \textbf{39} (2006)
1253--1255.

\bibitem{Wasow1965} W.~Wasow, \emph{Asymptotic Expansions for Ordinary Differential
Equations}, Wiley, 1965 (Dover repr.).

\bibitem{JMU1981} M.~Jimbo, T.~Miwa, K.~Ueno, \emph{Monodromy preserving deformation of
linear ODEs with rational coefficients. I}, Physica D \textbf{2} (1981) 306--352.

\bibitem{LR2016} M.~Loday-Richaud, \emph{Divergent Series, Summability and Resurgence II},
Lecture Notes in Math.\ \textbf{2154}, Springer, 2016.

\bibitem{Boalch2001} P.~Boalch, \emph{Symplectic manifolds and isomonodromic deformations},
Adv.\ Math.\ \textbf{163} (2001) 137--205.

\bibitem{Boalch2007} P.~Boalch, \emph{Quasi-Hamiltonian geometry of meromorphic connections},
Duke Math.\ J.\ \textbf{139} (2007) 369--405.

\bibitem{Boalch2014} P.~Boalch, \emph{Geometry and braiding of Stokes data; fission and wild
character varieties}, Ann.\ of Math.\ \textbf{179} (2014) 301--365, arXiv:1111.6228.

\bibitem{Katz1996} N.~M.~Katz, \emph{Rigid Local Systems}, Ann.\ of Math.\ Studies
\textbf{139}, Princeton, 1996.

\bibitem{DR2000} M.~Dettweiler, S.~Reiter, \emph{An algorithm of Katz and its application to
the inverse Galois problem}, J.\ Symbolic Comput.\ \textbf{30} (2000) 761--798.

\bibitem{DR2007} M.~Dettweiler, S.~Reiter, \emph{Middle convolution of Fuchsian systems and
the construction of rigid differential systems}, J.\ Algebra \textbf{318} (2007) 1--24.

\bibitem{Kolchin1973} E.~R.~Kolchin, \emph{Differential Algebra and Algebraic Groups},
Academic Press, 1973.

\bibitem{vdPS2003} M.~van der Put, M.~F.~Singer, \emph{Galois Theory of Linear Differential
Equations}, Grundlehren \textbf{328}, Springer, 2003.

\bibitem{Ramis1985} J.-P.~Ramis, \emph{Ph\'enom\`ene de Stokes et filtration Gevrey sur le
groupe de Picard--Vessiot}, C.\ R.\ Acad.\ Sci.\ Paris I \textbf{301} (1985) 165--167; full
proof in Martinet--Ramis, Ann.\ IHP \textbf{54} (1991) 331--401.

\bibitem{CS2007} P.~J.~Cassidy, M.~F.~Singer, \emph{Galois theory of parameterized
differential equations and linear differential algebraic groups}, IRMA Lect.\ Math.\ Theor.\
Phys.\ \textbf{9}, EMS, 2007, arXiv:math/0502396.

\bibitem{Nishioka1988} K.~Nishioka, \emph{A note on the transcendency of Painlev\'e's first
transcendent}, Nagoya Math.\ J.\ \textbf{109} (1988) 63--67.

\bibitem{Umemura1990} H.~Umemura, \emph{Second proof of the irreducibility of the first
differential equation of Painlev\'e}, Nagoya Math.\ J.\ \textbf{117} (1990) 125--171.

\bibitem{Casale2008} G.~Casale, \emph{Le groupo\"ide de Galois de $P_1$ et son
irr\'eductibilit\'e}, Comment.\ Math.\ Helv.\ \textbf{83} (2008) 471--519.

\bibitem{CL2009} S.~Cantat, F.~Loray, \emph{Dynamics on character varieties and Malgrange
irreducibility of Painlev\'e VI equation}, Ann.\ Inst.\ Fourier \textbf{59} (2009)
2927--2978.

\bibitem{Ronveaux1995} A.~Ronveaux (ed.), \emph{Heun's Differential Equations}, Oxford
Univ.\ Press, 1995.

\end{thebibliography}

\end{document}
